\subsection{Tree Recursion}
\label{Section 1.2.2}

Another common pattern of computation is called \newterm{tree recursion}.  As
an example, consider computing the sequence of Fibonacci numbers, in which each
number is the sum of the preceding two:
\begin{comment}
0, 1, 1, 2, 3, 5, 8, 13, 21, \( \dots \)
\end{comment}

$$ 0,\; 1,\; 1,\; 2,\; 3,\; 5,\; 8,\; 13,\; 21,\; \dots. $$

\noindent
In general, the Fibonacci numbers can be defined by the rule

$$ {\rm Fib}(n) =
\begin{cases}
        \; 0 & {\rm if} \;\; n=0, \\
	\; 1 & {\rm if} \;\; n=1, \\
	\; {\rm Fib}(n-1) + {\rm Fib}(n-2) \quad & {\rm otherwise}.
\end{cases} $$

\noindent
We can immediately translate this definition into a recursive procedure for
computing Fibonacci numbers:

\begin{codeBlock}
>>> def fib(n):
...     if n == 0:
...         return 0
...     elif n == 1:
...         return 1
...     else:
...         return fib(n - 1) + fib(n - 2)
\end{codeBlock}

\noindent
Consider the pattern of this computation.  To compute \code{fib(5)}, we
compute \code{fib(4)} and \code{fib(3)}.  To compute \code{fib(4)}, we
compute \code{fib(3)} and \code{fib(2)}.  In general, the evolved process
looks like a tree, as shown in \link{Figure 1.5}.  Notice that the branches
split into two at each level (except at the bottom); this reflects the fact
that the \code{fib} procedure calls itself twice each time it is invoked.

This procedure is instructive as a prototypical tree recursion, but it is a
terrible way to compute Fibonacci numbers because it does so much redundant
computation.  Notice in \link{Figure 1.5} that the entire computation of
\code{fib(3)}---almost half the work---is duplicated.  In fact, it is not hard
to show that the number of times the procedure will compute \code{fib(1)} or
\code{fib(0)} (the number of leaves in the above tree, in general) is
precisely Fib(\( n+1 \)).  To get an idea of how bad this is, one can
show that the value of Fib(\( n \)) grows exponentially with \( n \).  More
precisely (see \link{Exercise 1.13}), Fib(\( n \)) is the closest integer
to \( \varphi^n / \sqrt{5} \), where

$$\varphi = \frac{1 + \sqrt{5}}{2} \approx 1.6180 $$

\noindent
is the \newterm{golden ratio}, which satisfies the equation

$$\varphi^2 = \varphi + 1. $$


\begin{imageFigure}{The tree-recursive process generated in computing \code{fib(5)}}
  % The tree-recursive process generated in computing fib(5), used in Section 1.2.2.
%
% Upstream's figure also traces the order in which the process visits the
% nodes, with a dotted path threading the whole tree, and that is reproduced
% here: it shows that the process is depth-first -- the whole left subtree is
% computed before the right is begun -- which is what makes the space
% requirement the DEPTH of the tree rather than its size.
%
% Node positions are the leaves spaced evenly with every parent at the midpoint
% of its children, which is how upstream's layout is arranged too.

\begin{tikzpicture}[
  x = 12.2mm, y = -11mm,
  every node/.style = {font=\ttfamily\footnotesize, inner sep=1.2pt},
  edge/.style = {draw, thin},
]

% level 0
\node (f5)  at (4.3125, 0) {fib(5)};
% level 1
\node (f4)  at (2.375, 1) {fib(4)};
\node (f3b) at (6.25,  1) {fib(3)};
% level 2
\node (f3a) at (1.25,  2) {fib(3)};
\node (f2b) at (3.5,   2) {fib(2)};
\node (f2c) at (5.5,   2) {fib(2)};
\node (f1e) at (7,     2) {fib(1)};
% level 3
\node (f2a) at (0.5,   3) {fib(2)};
\node (f1b) at (2,     3) {fib(1)};
\node (f1c) at (3,     3) {fib(1)};
\node (f0b) at (4,     3) {fib(0)};
\node (f1d) at (5,     3) {fib(1)};
\node (f0c) at (6,     3) {fib(0)};
% level 4
\node (f1a) at (0,     4) {fib(1)};
\node (f0a) at (1,     4) {fib(0)};

% the value each terminal call returns, directly beneath it
\node (v1e) at (7,     3) {1};
\node (v1b) at (2,     4) {1};
\node (v1c) at (3,     4) {1};
\node (v0b) at (4,     4) {0};
\node (v1d) at (5,     4) {1};
\node (v0c) at (6,     4) {0};
\node (v1a) at (0,     5) {1};
\node (v0a) at (1,     5) {0};

\foreach \p/\c in {f5/f4, f5/f3b,
                   f4/f3a, f4/f2b, f3b/f2c, f3b/f1e,
                   f3a/f2a, f3a/f1b, f2b/f1c, f2b/f0b, f2c/f1d, f2c/f0c,
                   f2a/f1a, f2a/f0a}
  \draw[edge] (\p) to[out=-90, in=90] (\c);

\foreach \n/\v in {f1e/v1e, f1b/v1b, f1c/v1c, f0b/v0b,
                   f1d/v1d, f0c/v0c, f1a/v1a, f0a/v0a}
  \draw[edge] (\n) -- (\v);

% The order in which the process visits the nodes: down the left of each node,
% through its children, back up its right. Tracing this contour shows that the
% process is depth-first -- the whole left subtree is finished before the right
% is begun -- which is why the space it needs is the DEPTH of the tree and not
% its size, as the text goes on to say.
%
% Generated as an Euler tour of the tree rather than laid out by hand, so it
% cannot drift out of step with the nodes above.
\draw[dotted, thin, rounded corners=1.6pt, -{Triangle[open, length=1.4mm]}]
  (3.873,-0.44)
  -- (3.873,-0.44)
  -- (3.873,0.44)
  -- (1.935,0.56)
  -- (1.935,1.44)
  -- (0.81,1.56)
  -- (0.81,2.44)
  -- (0.06,2.56)
  -- (0.06,3.44)
  -- (-0.44,3.56)
  -- (-0.44,4.44)
  -- (-0.44,4.56)
  -- (-0.44,5.44)
  -- (0.44,5.44)
  -- (0.44,4.56)
  -- (0.44,4.44)
  -- (0.44,3.56)
  -- (0.56,3.56)
  -- (0.56,4.44)
  -- (0.56,4.56)
  -- (0.56,5.44)
  -- (1.44,5.44)
  -- (1.44,4.56)
  -- (1.44,4.44)
  -- (1.44,3.56)
  -- (0.94,3.44)
  -- (0.94,2.56)
  -- (1.56,2.56)
  -- (1.56,3.44)
  -- (1.56,3.56)
  -- (1.56,4.44)
  -- (2.44,4.44)
  -- (2.44,3.56)
  -- (2.44,3.44)
  -- (2.44,2.56)
  -- (1.69,2.44)
  -- (1.69,1.56)
  -- (3.06,1.56)
  -- (3.06,2.44)
  -- (2.56,2.56)
  -- (2.56,3.44)
  -- (2.56,3.56)
  -- (2.56,4.44)
  -- (3.44,4.44)
  -- (3.44,3.56)
  -- (3.44,3.44)
  -- (3.44,2.56)
  -- (3.56,2.56)
  -- (3.56,3.44)
  -- (3.56,3.56)
  -- (3.56,4.44)
  -- (4.44,4.44)
  -- (4.44,3.56)
  -- (4.44,3.44)
  -- (4.44,2.56)
  -- (3.94,2.44)
  -- (3.94,1.56)
  -- (2.815,1.44)
  -- (2.815,0.56)
  -- (5.81,0.56)
  -- (5.81,1.44)
  -- (5.06,1.56)
  -- (5.06,2.44)
  -- (4.56,2.56)
  -- (4.56,3.44)
  -- (4.56,3.56)
  -- (4.56,4.44)
  -- (5.44,4.44)
  -- (5.44,3.56)
  -- (5.44,3.44)
  -- (5.44,2.56)
  -- (5.56,2.56)
  -- (5.56,3.44)
  -- (5.56,3.56)
  -- (5.56,4.44)
  -- (6.44,4.44)
  -- (6.44,3.56)
  -- (6.44,3.44)
  -- (6.44,2.56)
  -- (5.94,2.44)
  -- (5.94,1.56)
  -- (6.56,1.56)
  -- (6.56,2.44)
  -- (6.56,2.56)
  -- (6.56,3.44)
  -- (7.44,3.44)
  -- (7.44,2.56)
  -- (7.44,2.44)
  -- (7.44,1.56)
  -- (6.69,1.44)
  -- (6.69,0.56)
  -- (4.753,0.44)
  -- (4.753,-0.44);

\end{tikzpicture}

\end{imageFigure}

\noindent
Thus, the process uses a number of steps that grows exponentially with the
input.  On the other hand, the space required grows only linearly with the
input, because we need keep track only of which nodes are above us in the tree
at any point in the computation.  In general, the number of steps required by a
tree-recursive process will be proportional to the number of nodes in the tree,
while the space required will be proportional to the maximum depth of the tree.

We can also formulate an iterative process for computing the Fibonacci numbers.
The idea is to use a pair of integers \( a \) and \( b \), initialized to
Fib(1) = 1 and Fib(0) = 0, and to repeatedly apply the
simultaneous transformations

$$
\begin{array}{l@{\;\;\gets\;\;}l}
  a & a + b, \\
  b & a.
\end{array}
$$

\noindent
It is not hard to show that, after applying this transformation \( n \) times,
\( a \) and \( b \) will be equal, respectively, to Fib(\( n+1 \)) and
Fib(\( n \)).  Thus, we can compute Fibonacci numbers iteratively using
the procedure

\begin{codeBlock}
>>> def fib(n):
...     return fib_iter(1, 0, n)

>>> @tail_recursive
... def fib_iter(a, b, count):
...     if count == 0:
...         return b
...     else:
...         return TailCall(fib_iter, a + b, a, count - 1)
\end{codeBlock}

\noindent
\begin{departure}
The original observes that the tree-recursive \code{fib} is inefficient, and
in Scheme that is the whole of the matter---\code{(fib 100)} is merely a
computation nobody will wait for.  In Python it is also \newterm{bounded}.  The
process descends \code{n} levels before the first value comes back, and Python
keeps a frame for every level, so somewhere around \code{fib(1000)} the
procedure stops being slow and starts being impossible.

That is a smaller distinction than it sounds, since nobody was going to wait
for \code{fib(100)} either.  But it is worth noticing that the two resources
the original is teaching us to count---steps and space---fail in different
ways here.  Running out of steps means waiting.  Running out of space means an
error.
\end{departure}

This second method for computing Fib(\( n \)) is a linear iteration.  The
difference in number of steps required by the two methods---one linear in
\( n \), one growing as fast as Fib(\( n \)) itself---is enormous, even for
small inputs.

One should not conclude from this that tree-recursive processes are useless.
When we consider processes that operate on hierarchically structured data
rather than numbers, we will find that tree recursion is a natural and powerful
tool.\footnote{An example of this was hinted at in \link{Section 1.1.3}. The
interpreter itself evaluates expressions using a tree-recursive process.} But
even in numerical operations, tree-recursive processes can be useful in helping
us to understand and design programs.  For instance, although the first
\code{fib} procedure is much less efficient than the second one, it is more
straightforward, being little more than a translation into Lisp of the
definition of the Fibonacci sequence.  To formulate the iterative algorithm
required noticing that the computation could be recast as an iteration with
three state variables.

\subsubsection*{Example: Counting change}

It takes only a bit of cleverness to come up with the iterative Fibonacci
algorithm.  In contrast, consider the following problem: How many different
ways can we make change of \$1.00, given half-dollars, quarters, dimes,
nickels, and pennies?  More generally, can we write a procedure to compute the
number of ways to change any given amount of money?

This problem has a simple solution as a recursive procedure.  Suppose we think
of the types of coins available as arranged in some order.  Then the following
relation holds:

\begin{niceQuote}
The number of ways to change amount \( a \) using \( n \) kinds of coins equals
\begin{itemize}
\item
the number of ways to change amount \( a \) using all but the first kind of coin,
plus
\item
the number of ways to change amount \( a - d \) using all \( n \) kinds of
coins, where \( d \) is the denomination of the first kind of coin.
\end{itemize}
\end{niceQuote}

\noindent
To see why this is true, observe that the ways to make change can be divided
into two groups: those that do not use any of the first kind of coin, and those
that do.  Therefore, the total number of ways to make change for some amount is
equal to the number of ways to make change for the amount without using any of
the first kind of coin, plus the number of ways to make change assuming that we
do use the first kind of coin.  But the latter number is equal to the number of
ways to make change for the amount that remains after using a coin of the first
kind.

Thus, we can recursively reduce the problem of changing a given amount to the
problem of changing smaller amounts using fewer kinds of coins.  Consider this
reduction rule carefully, and convince yourself that we can use it to describe
an algorithm if we specify the following degenerate cases:\footnote{For
example, work through in detail how the reduction rule applies to the problem
of making change for 10 cents using pennies and nickels.}

\begin{itemize}
\item
If \( a \) is exactly 0, we should count that as 1 way to make change.

\item
If \( a \) is less than 0, we should count that as 0 ways to make change.

\item
If \( n \) is 0, we should count that as 0 ways to make change.
\end{itemize}

\noindent
We can easily translate this description into a recursive procedure:

\begin{codeBlock}
>>> def count_change(amount):
...     return cc(amount, 5)

>>> def cc(amount, kinds_of_coins):
...     if amount == 0:
...         return 1
...     elif amount < 0 or kinds_of_coins == 0:
...         return 0
...     else:
...         return (cc(amount, kinds_of_coins - 1)
...                 + cc(amount - first_denomination(kinds_of_coins),
...                      kinds_of_coins))

>>> def first_denomination(kinds_of_coins):
...     if kinds_of_coins == 1:
...         return 1
...     elif kinds_of_coins == 2:
...         return 5
...     elif kinds_of_coins == 3:
...         return 10
...     elif kinds_of_coins == 4:
...         return 25
...     elif kinds_of_coins == 5:
...         return 50
\end{codeBlock}

\noindent
(The \code{first\_denomination} procedure takes as input the number of kinds of
coins available and returns the denomination of the first kind.  Here we are
thinking of the coins as arranged in order from largest to smallest, but any
order would do as well.)  We can now answer our original question about
changing a dollar:

\begin{codeBlock}
>>> count_change(100)
~\textit{292}~
\end{codeBlock}

\noindent
\code{count\_change} generates a tree-recursive process with redundancies
similar to those in our first implementation of \code{fib}.  (It will take
quite a while for that 292 to be computed.)  On the other hand, it is not
obvious how to design a better algorithm for computing the result, and we leave
this problem as a challenge.  The observation that a tree-recursive process may
be highly inefficient but often easy to specify and understand has led people
to propose that one could get the best of both worlds by designing a ``smart
compiler'' that could transform tree-recursive procedures into more efficient
procedures that compute the same result.\footnote{One approach to coping with
redundant computations is to arrange matters so that we automatically construct
a table of values as they are computed.  Each time we are asked to apply the
procedure to some argument, we first look to see if the value is already stored
in the table, in which case we avoid performing the redundant computation.
This strategy, known as \newterm{tabulation} or \newterm{memoization}, can be
implemented in a straightforward way.  Tabulation can sometimes be used to
transform processes that require an exponential number of steps (such as
\code{count\_change}) into processes whose space and time requirements grow
linearly with the input.  See \link{Exercise 3.27}.

A reader who knows Python will recognise this as \code{functools.cache}, which
is one line and needs no thought: written above \code{fib}, it takes
\code{fib(30)} from about 0.05 seconds to about 0.00002.  It is worth being
precise about what that fixes and what it does not.  The tree of redundant
calls is gone, so the number of \emph{steps} becomes linear.  The
\emph{depth} is untouched---the process still descends \code{n} levels before
it can return anything---so \code{fib(2000)} exhausts the stack whether the
values are remembered or not.  Memoization buys time at the price of space,
and space is the resource this process is already short of.  It also needs a
table that changes as the program runs, which is a larger idea than one line
suggests, and we will not use it until we have said what changing state costs.}

\begin{exerciseGroup}
\phantomsection\label{Exercise 1.11}
\exerciseHeading{Exercise 1.11:} A function \( f \) is defined by
the rule that

$$
f(n) =
\begin{cases}
\;\; n \quad \text{if \; \( n < 3 \),} \\
\;\; f(n-1) + 2\kern-0.08em f(n-2) + 3\kern-0.08em f(n-3) \quad \text{if \; \( n \ge 3 \).}
\end{cases}
$$

Write a procedure that computes \( f \) by means of a recursive process.  Write a procedure that
computes \( f \) by means of an iterative process.

\phantomsection\label{Exercise 1.12}
\exerciseHeading{Exercise 1.12:} The following pattern of numbers
is called \newterm{Pascal's triangle}.

\begin{example}
        1
      1   1
    1   2   1
  1   3   3   1
1   4   6   4   1
      . . .
\end{example}

The numbers at the edge of the triangle are all 1, and each number inside the
triangle is the sum of the two numbers above it.\footnote{The elements of
Pascal's triangle are called the \newterm{binomial coefficients}, because the
\( n^{\mathrm{th}} \) row consists of the coefficients of the terms in the expansion of
\( (x + y)^n \).  This pattern for computing the coefficients appeared in
Blaise Pascal's 1653 seminal work on probability theory, \textit{Trait\'e du
triangle arithm\'etique}.  According to \link{Knuth (1973)}, the same pattern appears
in the \textit{Szu-yuen Y\"u-chien} (``The Precious Mirror of the Four
Elements''), published by the Chinese mathematician Chu Shih-chieh in 1303, in
the works of the twelfth-century Persian poet and mathematician Omar Khayyam,
and in the works of the twelfth-century Hindu mathematician Bh\'ascara
\'Ach\'arya.} Write a procedure that computes elements of Pascal's triangle by
means of a recursive process.

\phantomsection\label{Exercise 1.13}
\exerciseHeading{Exercise 1.13:} Prove that Fib(\( n \)) is
the closest integer to \( \varphi^n / \sqrt{5} \), where \( \varphi = (1 +
\sqrt{5}) / 2 \).  Hint: Let \( \psi = (1 - \sqrt{5}) / 2 \).
Use induction and the definition of the Fibonacci numbers (see \link{Section 1.2.2})
to prove that \( \text{Fib}(n) = (\varphi^n - \psi^n) / \sqrt{5} \).
\end{exerciseGroup}

